\section{The massive Fredholm index}
\label{sec:massive-index}

We compute the index on the fixed Hilbert space
\[
 \mathscr H_G=\bigl[L^2(\mathfrak k^9)\otimes\mathcal F_N\bigr]^{\U(N)\times G},
 \qquad
 \mathscr H_G=\mathscr H_G^+\oplus\mathscr H_G^-.
\]
Here the superscripts denote the exterior grading fixed in
Section~\ref{sec:algebra}. Let \(D_h=D_u(h,1)|_{\mathscr H_G}\).
It is odd and self-adjoint by Lemma~\ref{lem:domains}.

\begin{samepage}
\begin{proposition}[Massive index]
\label{prop:massive-index}
For every \(h>0\), the graph domain of \(D_h\) embeds compactly into
\(\mathscr H_G\), and its even-to-odd part is Fredholm. Its index is
\begin{equation}
\label{eq:massive-index}
 \operatorname{ind}_G D_u(h,1)
 :=\dim\Ker(D_h|_{\mathscr H_G^+})
   -\dim\Ker(D_h|_{\mathscr H_G^-})=p(N),
\end{equation}
where \(p(N)\) is the number of unordered integer partitions of
\(N\). For all sufficiently small \(h>0\), the massive kernel itself
has dimension \(p(N)\), consists of even vectors, and is separated
from the rest of the spectrum of \(D_h\) by a positive gap uniform
in \(h\).
\end{proposition}
\end{samepage}

The proof occupies this section. It uses the massive deformation to
obtain a Fredholm problem, but makes no assertion yet about its
relation to the undeformed kernel. That relation will require the
global estimates of the subsequent sections.

\subsection{Confinement and continuity at positive coupling}

\begin{lemma}[Coercivity of the massive potential]
\label{lem:massive-potential}
There are \(c>0\) and \(R<\infty\), depending only on \(N\), such
that
\begin{equation}
\label{mass:potential-infinity}
 U_{1,1}(x)\geq c\abs{x}^2\qquad(\abs{x}\geq R).
\end{equation}
Moreover, for \(h,m>0\),
\begin{equation}
\label{mass:potential-scaling}
 U_{h,m}(x)=\frac{m^4}{h^2}
 U_{1,1}\left(\frac hm x\right).
\end{equation}
\end{lemma}

\begin{proof}
The scaling identity follows term by term from
\eqref{eq:square}. To prove the lower bound, suppose otherwise and
choose \(x^{(n)}\) with \(r_n=\abs{x^{(n)}}\to\infty\) and
\(U_{1,1}(x^{(n)})/r_n^2\to0\). After passing to a subsequence,
\(x_i^{(n)}/r_n\to a_i\). The last-six mass terms give \(a_d=0\),
while the first-three squares give
\begin{equation}
\label{mass:escaping-residual}
 [x_a^{(n)},x_b^{(n)}]-2x_c^{(n)}=o(r_n)
 \quad\text{for cyclic }(a,b,c).
\end{equation}
Dividing by \(r_n^2\) shows that the \(a_a\) commute. Their joint
eigenspaces give an orthogonal decomposition
\(\C^N=\bigoplus_\alpha E_\alpha\), with distinct joint eigenvalue
triples \(\lambda_\alpha\in\R^3\). Choose \(t\in\R^3\) so that
the numbers \(t\cdot\lambda_\alpha\) are distinct, and put
\(y_n=\sum_a t_ax_a^{(n)}\).

The spectral subspaces of \(y_n\) in the corresponding clusters
converge to the \(E_\alpha\). By unitaries tending to the identity
we may identify these subspaces with fixed blocks, without changing
the preceding limits. Their spectral gaps are at least \(c_0r_n\)
for some \(c_0>0\). Equation~\eqref{mass:escaping-residual} implies
\(\norm{[y_n,x_a^{(n)}]}=O(r_n)\). In bases diagonalizing the two
relevant blocks of \(y_n\), an off-block matrix entry is multiplied
in this commutator by a spectral difference of magnitude at least
\(c_0r_n\). Hence
\begin{equation}
\label{mass:off-block-bound}
 \norm{(x_a^{(n)})_{\alpha\beta}}=O(1),\qquad\alpha\ne\beta.
\end{equation}
This is a finite-dimensional Hilbert--Schmidt estimate, uniform in
the sequence.

The trace in a diagonal block of \([x_a^{(n)},x_b^{(n)}]\) has zero
contribution from its internal commutator. Its remaining terms are
products of the off-block entries in \eqref{mass:off-block-bound},
so its block trace is \(O(1)\). Divide the block trace of
\eqref{mass:escaping-residual} by \(r_n\). The result is
\(2\dim(E_\alpha)(\lambda_\alpha)_c=0\) for every \(\alpha,c\).
Thus every \(a_i\) vanishes. This contradicts
\(\sum_i\abs{a_i}^2=1\), proving \eqref{mass:potential-infinity}.
Grouping equal limiting joint eigenvalues into a single block was
essential here; no gap was required inside such a block.
\end{proof}

\begin{samepage}
\begin{lemma}[A common massive graph domain]
\label{mass:common-domain}
The domain of \(D_h\), for every \(h>0\), is
\begin{equation}
\label{mass:graph-domain}
 \mathcal D=
 \left\{\psi\in\mathscr H_G\cap H^1:
  \abs{x}\psi\in L^2,\quad [x_i,x_j]\psi\in L^2\ (i<j)\right\}.
\end{equation}
On each compact interval \(I\Subset(0,\infty)\), its graph norm is
uniformly equivalent to
\begin{equation}
\label{mass:graph-norm}
 \norm{\psi}+\norm{\nabla\psi}+\norm{\abs{x}\psi}
       +\sum_{i<j}\norm{[x_i,x_j]\psi}.
\end{equation}
The embedding \(\mathcal D\hookrightarrow\mathscr H_G\) is compact,
and \(h\mapsto D_h\) is norm-resolvent continuous on \((0,\infty)\).
Its graded Fredholm index is constant there.
\end{lemma}
\end{samepage}

\begin{proof}
The fermionic multiplier in \eqref{eq:square}, denoted below by
\(V_F(h,m;x)\), acts on a fixed finite-dimensional space and satisfies
\begin{equation}
\label{mass:fermion-bound}
 \norm{V_F(h,m;x)}_{\mathrm{op}}
 \leq C_N(m+h\abs{x}).
\end{equation}
Lemma~\ref{lem:massive-potential} and scaling give, uniformly for
\(h\in I\), \(U_{h,1}\geq c_I\abs{x}^2-C_I\). Absorbing the
linear term in \eqref{mass:fermion-bound} yields
\[
 \cH(h,1)[\psi]+C_I\norm{\psi}^2
 \geq \frac12\norm{\nabla\psi}^2+c'_I\norm{\abs{x}\psi}^2.
\]
The identity for \(\cH\) then also bounds
\(\int U_{h,1}\abs{\psi}^2\) by
\(C_I(\cH(h,1)[\psi]+\norm{\psi}^2)\). Each bracket multiplier
is controlled by this integral and \(\norm{\abs{x}\psi}\), because
the first-three residuals are \(h[x_a,x_b]-2x_c\) and all the other
brackets occur directly in the potential. On \(\mathscr H_G\),
Lemma~\ref{lem:coercive-charge} consequently bounds
\eqref{mass:graph-norm} by \(C_I(\norm{D_h\psi}+\norm{\psi})\).
The reverse bound follows immediately from the formula for \(D_h\).

These statements first hold on the invariant core. Its graph
completion equals the maximal domain by Lemma~\ref{lem:domains}.
Conversely a vector satisfying \eqref{mass:graph-domain} has all
terms of the charge in \(L^2\), so is in that maximal domain.
This proves the common-domain assertion and norm equivalence.
For a bounded set in this graph norm, Rellich compactness on balls
gives local \(L^2\) compactness, whereas
\[
 \norm{\mathbf1_{\{\abs{x}>R\}}\psi}^2
 \leq R^{-2}\norm{\abs{x}\psi}^2
\]
makes the tails uniformly small. The graph embedding is compact.
Thus \(D_h\) has compact resolvent, finite-dimensional kernel, and
closed range, and its odd part is Fredholm.

Writing \(D_h-D_{h_0}=(h-h_0)B\), the bracket bound shows that
\(B:\mathcal D\to\mathscr H_G\) is bounded for the graph norm of
\(D_{h_0}\). Thus
\((D_h-D_{h_0})(D_{h_0}-\ii)^{-1}\) tends to zero in operator norm.
The resolvent identity and a Neumann series imply norm-resolvent
continuity. For completeness, choose a symmetric interval
\((-a,a)\) whose endpoints avoid the spectrum at \(h_0\). Its
finite-rank spectral projection is norm-continuous for nearby
\(h\), and preserves parity. The even and odd ranks of the
projection are therefore locally constant. This is the standard stability
of separated spectral projections; see \cite[Chapter~IV, Section~3.4,
Theorem~3.16]{Kato1995}. The graph estimates above and the parity argument
here are specific to the present operators. On its nonzero spectral
subspace, the odd operator \(D_h\) is an isomorphism between the
even and odd parts, so their rank difference is exactly the graded
kernel dimension. The index is locally constant, and hence
constant on the connected interval \((0,\infty)\).
\end{proof}

\subsection{The minima and their normal oscillators}

The zero set \(Z=\{U_{1,1}=0\}\) is a finite union of compact gauge
orbits. In fact its equations are
\[
 x_a=2j_a,\qquad [j_a,j_b]=\epsilon_{abc}j_c,\qquad x_d=0.
\]
Thus \(j\) is a unitary representation of \(\mathfrak{su}(2)\) on
\(\C^N\), and its irreducible dimensions form a partition of \(N\).
Conversely every such representation gives a zero. Conjugacy by
\(\U(N)\) is the same as conjugacy by \(\SU(N)\) on this
configuration space, and the representing matrices have trace zero.
There are therefore exactly \(p(N)\) zero orbits.
These partition labels are the classical vacuum labels of the plane-wave
model~\cite[Sec.~5]{BMN2002}.

For each zero orbit, the required contribution is one even normal
vacuum that survives the stabilizer symmetry. We first specify the
normal Hamiltonian, then calculate its ground line and its character.

Fix a representative \(q=(2j_1,2j_2,2j_3,0,\ldots,0)\in Z\).
Let \(\mathcal O=\U(N)q\), let \(H_q\) be its stabilizer, and set
\(S_q=(T_q\mathcal O)^\perp\subset\mathfrak k^9\). Write
\(U_q^{(2)}\) for the quadratic Taylor term of \(U_{1,1}\) at
\(q\), restricted to \(S_q\). The limiting normal Hamiltonian is
\begin{equation}
\label{mass:normal-Hamiltonian}
 H_q^{\mathrm{nor}}
 =-\frac12\Delta_{S_q}+U_q^{(2)}(z)
   +\frac{\ii}{2}\theta^t
      \left(\frac32T-2\ad(j_a)\gamma^a\right)\theta.
\end{equation}
It acts first on the full slice space
\(L^2(S_q)\otimes\mathcal F_N\); gauge transport will impose
\(H_q\)-equivariance there.

The fluctuation frequencies below were computed by
Dasgupta--Sheikh-Jabbari--Van Raamsdonk
\cite[Sections~5.2--5.4]{DasguptaSheikhJabbariVanRaamsdonk2002}.
We rederive them through spin coupling to specify the normal Hessian
and then establish the stabilizer character needed for the invariant
index. In a finite
unitary \(\mathfrak{su}(2)\)-representation, the Hermitian third
generator has a highest weight \(l\). The adjoint raising and
lowering operators produce weights \(l,l-1,\ldots,-l\); on a unit
weight-\(r\) vector the squared lowering norm is
\(l(l+1)-r(r-1)\). This follows by moving a raising operator past a
lowering operator using their commutator. Positivity terminates the
string at \(-l\), so \(2l\) is a nonnegative integer. These strings
split off orthogonally and have uniquely determined ladder
coefficients. Tensoring two strings and adding their weights gives
the usual decomposition into consecutive spins
\(\abs{l-s},\ldots,l+s\). For real skew generators, whose squares
have Casimir \(-l(l+1)\), this implies
\begin{equation}
\label{mass:spin-coupling}
 \sum_a\ad(j_a)L_a^{(s)}
 =\frac{l(l+1)+s(s+1)-k(k+1)}2
 \quad\text{on total spin }k.
\end{equation}

\begin{lemma}[Normal frequencies and the vacuum line]
\label{lem:normal-oscillators}
The quadratic form \(U_q^{(2)}\) is positive definite on \(S_q\).
The oscillator \(H_q^{\mathrm{nor}}\) is nonnegative, has a unique
zero-energy line, and has a positive spectral gap on its orthogonal
complement. Its zero-energy line is even and is fixed by the entire
stabilizer \(H_q\), including exchanges of identical irreducible
summands. It is also fixed by the combined first-three spatial and
compensating gauge rotations, and by \(\Spin(6)\).
\end{lemma}

\begin{proof}
\emph{Bosonic normal modes.}
Decompose \(\mathfrak k_\C\) into irreducible spin-\(l\) summands
for \(\ad(j)\). Half-integral spins are allowed: maps between two
different irreducible summands of the defining representation may
carry such spins. The dimensions calculated in the complexification
give the real dimensions of the original real operators after
summing all their summands.

On the first three spatial components, the linearized residual in
the potential is
\[
 2\left(\sum_a\ad(j_a)L_a^{(1)}-\Id\right),
\]
where \(L_a^{(1)}\) generates cross product on \(\R^3\).
Equation~\eqref{mass:spin-coupling} gives the eigenvalues
\(-l,1,l+1\) of \(\sum_a\ad(j_a)L_a^{(1)}\), on coupled spins
\(l+1,l,l-1\), respectively, whenever the indicated representation
exists. Taking the absolute values of the linearized residual gives
normal frequencies
\begin{equation}
\label{mass:boson-frequencies}
 \begin{array}{c|c}
  \text{frequency}&\text{multiplicity}\\ \hline
  2(l+1)&2l+3\\
  2l&\max(2l-1,0)\\
  2l+1&6(2l+1).
 \end{array}
\end{equation}
The last row comes from the last six components, whose squared
frequency is \(1+4l(l+1)=(2l+1)^2\). For \(l>0\), the omitted
zero-frequency spin-\(l\) summand is exactly the gauge tangent:
\(w\mapsto([w,j_a])_a\) is an injective intertwiner from that
color summand, its image is tangent to the zero orbit, and the
dimensions agree. For \(l=0\), only the spin-one first-three
summand exists and has frequency two; there is no gauge tangent.
Thus all frequencies on \(S_q\) are positive. Their bosonic ground
energy per color spin-\(l\) summand is
\begin{equation}
\label{mass:boson-energy}
 E_B(l)=\frac12\left((2l+2)(2l+3)
              +(2l)(2l-1)+6(2l+1)^2\right)
        =16l^2+16l+6.
\end{equation}
The middle term is omitted when its representation is absent; it
is already zero for \(l=0\) and \(l=1/2\).

\smallskip\noindent\emph{Fermionic energy and the ground line.}
The real skew spin-one-half matrices
\(L_a^{(1/2)}=-T\gamma^a/2\) satisfy the same
\(\mathfrak{su}(2)\) relations, and the fermion frequency matrix is
\[
 M_q=\frac32T-2\ad(j_a)\gamma^a
   =T\left(\frac32-4\sum_a\ad(j_a)L_a^{(1/2)}\right).
\]
The coupled spins \(l+1/2\) and \(l-1/2\) have scalar factors
\(2l+3/2\) and \(-2l-1/2\), respectively, multiplying \(T\).
Since the complexified real sixteen-dimensional spinor has eight
copies of spin one-half, the absolute frequencies and their
complexified multiplicities are
\begin{equation}
\label{mass:fermion-frequencies}
 \begin{array}{c|c}
 \text{absolute frequency}&\text{complexified multiplicity}\\ \hline
 2l+3/2&8(2l+2)\\
 2l+1/2&8(2l).
 \end{array}
\end{equation}
There are no zero fermion frequencies. A real two-plane with skew
frequency \(\omega\) contributes eigenvalues \(\pm\omega/2\) to
\(\ii\theta^tM_q\theta/2\), by the CAR normalization. The fermionic
ground energy is consequently \(-\Tr\abs{M_q}/4\). From
\eqref{mass:fermion-frequencies} this is \(-E_B(l)\), with
\(E_B(l)\) as in \eqref{mass:boson-energy}. Every bosonic oscillator
and every fermionic two-plane has a unique ground vector, so their
product is the unique zero-energy vector and has a positive gap
above it.

\smallskip\noindent\emph{Stabilizer action and parity.}
The unique ground line contributes to the invariant index only if it
has trivial stabilizer character. We verify that condition in the
fixed exterior realization, which also determines its parity.
Decompose the defining representation as
\[
 \C^N=\bigoplus_r(V_{j_r}\otimes M_r),\qquad
 H_q=\prod_r\U(M_r),\qquad m_r=\dim M_r,
\]
where the \(j_r\) are distinct and \(V_j\) denotes spin \(j\).
The empty vector in the \(T\)-polarized Fock space is gauge
invariant. The one-particle spin representation under
\(\Spin(3)\times\Spin(6)\) is \(\C^2\otimes\C^4\). Relative to
this polarization, the fermionic ground vector fills precisely the
negative-factor space of \(M_q/T\). On a color spin-\(l\) summand
this is \(V_{l-1/2}\otimes\C^4\), of complex dimension
\begin{equation}
\label{mass:occupied-dimension}
 q_l=4(2l)=8l\qquad(l>0).
\end{equation}
Notice the distinction from the doubled multiplicities of the
complexified real frequency matrix in
\eqref{mass:fermion-frequencies}.

The color summand from maps between defining isotypic components
is
\[
 \operatorname{Hom}(V_{j_r},V_{j_s})
       \otimes\operatorname{Hom}(M_r,M_s),
\]
and its spin list is \(l=\abs{j_r-j_s},\ldots,j_r+j_s\), once
each. Thus the determinant character of its filled space is the
product over that list, omitting \(l=0\), of
\begin{equation}
\label{mass:centralizer-character}
 \left((\det g_s)^{m_r}(\det g_r)^{-m_s}\right)^{q_l}.
\end{equation}
The reverse pair \((s,r)\) has the identical spin list and the
inverse character, so the factors cancel. For \(r=s\) the displayed
character is already trivial. Removing the scalar color from
\(\operatorname{End}(\C^N)\) removes only an \(l=0\) summand and
does not affect this argument. Hence the determinant of the full
filled space is the trivial character of \(H_q\). The bosonic
ground Gaussian is also \(H_q\)-invariant. The entire ground line
is therefore fixed by the stabilizer, including its permutation
matrices exchanging identical summands. No ordering sign remains.
Moreover every \(q_l\) is a multiple of four, including when
\(l\) is half-integral. The occupied degree is even, proving the
parity assertion.

\smallskip\noindent\emph{Spatial rotations.}
First-three rotations of \(q\) can be compensated by the
\(\SU(2)\) representation defined by \(j\). Their combined action,
and the action of \(\Spin(6)\), preserve \(S_q\) and commute with
\(H_q^{\mathrm{nor}}\). Each acts on its unique ground line by a
continuous character of a connected semisimple group. Such a
character is trivial: its derivative annihilates the commutator
Lie algebra, which is the entire Lie algebra, and connectedness
then annihilates the character. This proves all the claimed
symmetries.
\end{proof}

\subsection{Small-coupling localization with norm exhaustion}

We next justify the passage from the finitely many normal
oscillators to the spectrum. This is harmonic comparison by
localization, as in semiclassical analysis~\cite{Simon1983}; here the
wells are gauge orbits, so their slice measures and stabilizers
must be retained explicitly. In particular, the argument includes
tails away from their tubes; local oscillator convergence alone
would not determine an index.

\begin{lemma}[Uniform localization near the massive zero orbits]
\label{mass:distance-confinement}
For \(0<h\leq1\), let
\(d_h(x)=\dist(x,h^{-1}Z)\). There are constants \(c,C>0\),
independent of \(h\), such that
\begin{equation}
\label{mass:distance-bound}
 U_{h,1}(x)\geq c\,d_h(x)^2,
 \qquad
 \cH(h,1)[\psi]\geq
 \frac12\norm{\nabla\psi}^2
 +\frac c2\norm{d_h\psi}^2-C\norm{\psi}^2.
\end{equation}
In particular, bounded form-energy families with bounded norm have
uniform kinetic bounds and
\(\norm{\mathbf1_{\{d_h>R\}}\psi}^2\leq C'/R^2\).
\end{lemma}

\begin{proof}
The positive normal Hessians in
Lemma~\ref{lem:normal-oscillators}, and the finitely many compact
orbits, imply
\(U_{1,1}\geq c\dist(\cdot,Z)^2\) in a fixed neighborhood of
\(Z\). Indeed normal tubular coordinates have positive quadratic
Taylor term and uniformly cubic Taylor remainder on these compact
orbits. On a compact set outside that neighborhood the same
inequality follows from the positive minimum of
\(U_{1,1}/\dist(\cdot,Z)^2\). At infinity it follows from
Lemma~\ref{lem:massive-potential}, since \(Z\) is bounded. Scaling
then gives its first assertion for all \(x\).

Let \(M=\max_{q\in Z}\abs{q}\). A nearest point of \(h^{-1}Z\)
gives \(h\abs{x}\leq hd_h(x)+M\). Equation~\eqref{mass:fermion-bound}
therefore implies
\[
 U_{h,1}(x)+V_F(h,1;x)
 \geq c\,d_h(x)^2-C_N(1+M+hd_h(x))
 \geq \frac c2d_h(x)^2-C.
\]
Adding the kinetic term proves \eqref{mass:distance-bound}.
In particular the matrix-valued potential plus fermion term is
nonnegative whenever \(d_h\geq R_0\), for a fixed \(R_0\), and
the claimed tail bound follows by integration.
\end{proof}

\begin{lemma}[Normal-slice compactness]
\label{mass:slice-compactness}
Suppose \(h_n\downarrow0\) and \(\psi_n\) are gauge-invariant
vectors of bounded norm and bounded \(\cH(h_n,1)\) form energy.
After a subsequence, there is one profile
\(f_q\in L^2(S_q)\otimes\mathcal F_N\) for each of the \(p(N)\) zero
orbits, with the following properties. Each profile is
\(H_q\)-equivariant and belongs to the oscillator form domain.
With the normalizations described in the proof,
\begin{equation}
\label{mass:slice-exhaustion}
 \lim_{n\to\infty}\norm{\psi_n}^2
   =\sum_q\norm{f_q}^2,
 \qquad
 \sum_q H_q^{\mathrm{nor}}[f_q]
   \leq\liminf_{n\to\infty}\cH(h_n,1)[\psi_n],
\end{equation}
whenever the subsequence realizes the norm limit and energy
liminf. For any fixed finite family of such sequences, extraction
can be simultaneous and the norm identity polarizes to their inner
products. Parity is preserved.
\end{lemma}

\begin{proof}
Choose disjoint normal tubes of radius \(\rho\) about the compact
orbits in \(Z\). The tubular neighborhood theorem here applies to
the smooth compact embedded orbit \(\U(N)/H_q\) in the Euclidean
configuration space, with its ordinary normal bundle. Scaling its
normal exponential map gives the exact coordinates
\begin{equation}
\label{mass:tube-chart}
 \U(N)\times_{H_q}B_{S_q}(\rho/h)
 \longrightarrow\mathfrak k^9,
 \qquad [k,z]\longmapsto\operatorname{Ad}(k)(q/h+z).
\end{equation}
Within this tube, \(d_h=\abs{z}\). Stabilizer equivariance is kept
on the linear slice: if \(\rho_F\) denotes the gauge action on
\(\mathcal F_N\), then restriction to the slice satisfies
\[
 \psi(q/h+\operatorname{Ad}(a)z)=\rho_F(a)\psi(q/h+z),
 \qquad a\in H_q.
\]
We do not replace this equivariant linear slice by a singular
quotient.

Let \(d_q=\dim\mathcal O\), and give \(\U(N)/H_q\) a fixed
normalized invariant measure. Integration in
\eqref{mass:tube-chart} gives measure
\(h^{-d_q}J_q(hz)\,\dd z\) after orbit integration, with
\(J_q\) smooth, positive, and \(H_q\)-invariant on
\(B_{S_q}(\rho)\). This follows by differentiating the tube map:
the \(d_q\) orbit columns are
\(h^{-1}[w,q+hz]\), while the normal columns are fixed orthonormal
vectors. Put
\(j_{q,h}(z)=J_q(hz)/J_q(0)\). On every bounded slice set,
\(j_{q,h}\to1\) smoothly and
\(\nabla\log j_{q,h}=O(h)\). Define the flat-measure slice vector
\begin{equation}
\label{mass:slice-normalization}
 f_{q,n}(z)=
 h_n^{-d_q/2}J_q(0)^{1/2}j_{q,h_n}(z)^{1/2}
 \psi_n(q/h_n+z),\qquad \abs{z}<\rho/h_n.
\end{equation}
This identification is isometric on the tube. For general form-domain
vectors we define it by \(L^2\) closure from smooth equivariant
vectors; the derivative estimates below pass by form-core
approximation. Distinct orbit dimensions cause no ambiguity: each
tube has its own displayed constant normalization.

For any smooth equivariant vector, differentiation in an orthonormal
basis of the fixed space \(S_q\) is ordinary ambient
differentiation. Pointwise its squared slice derivative is therefore
at most the full squared ambient gradient. After
\eqref{mass:slice-normalization}, the density derivative changes this
bound only by \(O(h_n)\) on bounded sets. The uniform kinetic bound
in Lemma~\ref{mass:distance-confinement} yields local \(H^1\) bounds
for every \(f_{q,n}\). Rellich compactness, followed by a diagonal
subsequence on expanding balls in the finitely many slices, gives
strong local \(L^2\) convergence and weak local \(H^1\) convergence
to \(f_q\). The distance bound gives both
\(\abs{z}f_q\in L^2\) and uniformly small squared norm tails
\(O(R^{-2})\). For fixed \(R\), all the tubes \(\abs{z}<R\)
are disjoint when \(n\) is large, and their complement is contained
in \(\{d_{h_n}\geq R\}\). Strong local convergence followed by
\(R\to\infty\) proves the first identity in
\eqref{mass:slice-exhaustion}. The same estimates for a finite list,
or the Cauchy--Schwarz inequality on its common tails, prove the
polarized identity. All changes of coordinates and densities commute
with parity, so parity passes to the profiles.

On a fixed bounded slice, Taylor expansion gives, uniformly there,
\begin{equation}
\label{mass:slice-potential-limit}
 U_{h,1}(q/h+z)=U_q^{(2)}(z)+O_R(h),\qquad
 V_F(h,1;q/h+z)=\frac{\ii}{2}\theta^tM_q\theta+O_R(h).
\end{equation}
For the first formula use
\(U_{h,1}(q/h+z)=h^{-2}U_{1,1}(q+hz)\) and the vanishing of the
constant and linear Taylor terms; for the second use the explicit
linear fermion multiplier. Thus local potential energies converge
under strong \(L^2\) convergence, and the slice-gradient inequality
and weak lower semicontinuity give the kinetic liminf.

To justify the global energy statement, take real gauge-invariant
cutoffs \(\chi_{q,h,R}\), equal to one in the \(R\)-tube about
\(h^{-1}\mathcal O_q\) and supported in its \(2R\)-tube. For
fixed \(R\) and small \(h\) the tubes are disjoint. Complete them
to a squared partition of unity with an exterior cutoff, choosing
the transition profiles so that the sum of squared gradients is
at most \(C/R^2\). The first-order symbol identity
\eqref{alg:symbol} gives the localization formula
\begin{equation}
\label{mass:IMS}
 \sum_\alpha\cH(h,1)[\chi_\alpha\psi]
 =\cH(h,1)[\psi]
    +\frac12\int\sum_\alpha\abs{\nabla\chi_\alpha}^2\abs{\psi}^2.
\end{equation}
Every localized gauge-invariant form on the left is nonnegative,
so the exterior term may be discarded. On each remaining tube the
cutoff is a fixed function \(\chi_R(z)\). Applying the preceding
local liminf yields
\[
 \liminf_n\cH(h_n,1)[\psi_n]
 \geq\sum_qH_q^{\mathrm{nor}}[\chi_Rf_q]-C/R^2.
\]
The uniform gradient and second-moment bounds already show that
the limits belong to the confining oscillator form domains.
Accordingly \(\chi_Rf_q\to f_q\) in these form domains as
\(R\to\infty\). This proves the second assertion of
\eqref{mass:slice-exhaustion} and completes the proof.
\end{proof}

\begin{lemma}[Recovery of each normal vacuum]
\label{mass:vacuum-recovery}
For every zero orbit there is a family \(\Phi_{q,h}\in\mathcal D\)
of even gauge- and \(\Spin(3)\times\Spin(6)\)-invariant vectors
such that, as \(h\downarrow0\),
\begin{equation}
\label{mass:vacuum-recovery-eq}
 \ip{\Phi_{q,h}}{\Phi_{q',h}}\longrightarrow\delta_{qq'},
 \qquad \norm{D_h\Phi_{q,h}}\longrightarrow0.
\end{equation}
\end{lemma}

\begin{proof}
Let \(\varphi_q\) be a normalized ground vector of
\(H_q^{\mathrm{nor}}\), as in
Lemma~\ref{lem:normal-oscillators}, and first fix a radial slice
cutoff \(\chi_R\) supported in \(\abs{z}<2R\). Use
\eqref{mass:slice-normalization} inversely to put
\(\chi_R\varphi_q\) in the tube, and extend by gauge equivariance
and then by zero. The stabilizer character calculation makes this
a well-defined vector, denoted \(\Phi_{q,h,R}\). Its support is
inside the tube for all sufficiently small \(h\), and it is in
the charge domain. The density and radial cutoff are preserved by
the combined spatial and compensating gauge rotations. The
invariance in Lemma~\ref{lem:normal-oscillators} therefore makes
\(\Phi_{q,h,R}\) invariant under the full deformation rotation
group. It is even as well.

We verify energy convergence, including the derivative directions
not contained in the fixed normal space. On a bounded slice the
orbit columns in \eqref{mass:tube-chart} are
\([w,q/h+z]\), with \(w\) in a fixed complement of the
stabilizer Lie algebra. At \(z=0\) they span the orthogonal
complement of \(S_q\), with inverse of size \(O(h)\). On
\(\abs{z}\leq2R\) their inverse is still \(O_R(h)\). Gauge
equivariance differentiates an orbit column to the fixed
finite-dimensional operator \(\rho_{F,*}(w)\). Expressing a unit
ambient tangent vector in these columns and the normal columns
therefore bounds the additional derivatives by
\[
 C_Rh\bigl(\abs{f}+\abs{z}\abs{\nabla f}\bigr)
 \quad\text{for the corresponding fixed slice field }f.
\]
The density derivative is also \(O_R(h)\). It follows that the
kinetic form of \(\Phi_{q,h,R}\) converges to the Euclidean
slice kinetic form of \(\chi_R\varphi_q\). Together with
\eqref{mass:slice-potential-limit}, this gives
\[
 \norm{\Phi_{q,h,R}}\longrightarrow\norm{\chi_R\varphi_q},\qquad
 \cH(h,1)[\Phi_{q,h,R}]
 \longrightarrow H_q^{\mathrm{nor}}[\chi_R\varphi_q].
\]
The latter is
\(\frac12\norm{(\nabla\chi_R)\varphi_q}^2\), since
\(H_q^{\mathrm{nor}}\varphi_q=0\); it tends to zero as
\(R\to\infty\). Rotation invariance and
Proposition~\ref{prop:square} identify this massive form energy
with \(\norm{D_h\Phi_{q,h,R}}^2\). Choose \(R\to\infty\)
sufficiently slowly as \(h\downarrow0\) that all preceding
fixed-\(R\) convergences hold. This is a diagonal choice: for each
integer \(R\), first choose a positive upper bound on \(h\)
making its finitely many errors smaller than \(1/R\). Supports for
different orbits are disjoint for these choices. Normalizing the
vectors proves \eqref{mass:vacuum-recovery-eq}.
\end{proof}

\subsection{From localization to the index}

\begin{proof}[Proof of Proposition~\ref{prop:massive-index}]
The compact-domain and continuity assertions were proved in
Lemma~\ref{mass:common-domain}. We compute the constant index as
\(h\downarrow0\).

There exist \(\delta>0\) and \(h_0>0\) such that
\begin{equation}
\label{mass:odd-gap}
 \norm{D_h\psi}\geq\delta\norm{\psi},
 \qquad \psi\in\mathcal D\cap\mathscr H_G^-,\quad 0<h<h_0.
\end{equation}
Otherwise one could choose normalized odd \(\psi_n\), with
\(h_n\downarrow0\) and \(\norm{D_{h_n}\psi_n}\to0\).
Lemma~\ref{lem:coercive-charge} would give
\(\cH(h_n,1)[\psi_n]\to0\). By
Lemma~\ref{mass:slice-compactness}, their profiles would have total
squared norm one and would all belong to the zero-energy lines of
the normal oscillators. Every such line is even by
Lemma~\ref{lem:normal-oscillators}, whereas every extracted profile
is odd. This contradiction proves \eqref{mass:odd-gap}.

Because \(D_h\) is odd and self-adjoint on the same invariant
Hilbert space, it pairs positive even and odd eigenvectors of
\(D_h^2\): for eigenvalue \(\lambda>0\), the map
\(\lambda^{-1/2}D_h\) is an isometric parity-changing bijection.
Compact resolvent makes this a statement about a complete discrete
spectrum. Thus \eqref{mass:odd-gap} excludes positive even
eigenvalues below \(\delta^2\) as well, and the kernel is entirely
even. If \(P_h\) denotes its orthogonal projection, we obtain
\begin{equation}
\label{mass:kernel-gap}
 \norm{(1-P_h)\psi}\leq\delta^{-1}\norm{D_h\psi},
 \qquad 0<h<h_0.
\end{equation}

The kernel has dimension at most \(p(N)\) for all sufficiently
small \(h\). If not, along a sequence \(h_n\downarrow0\) one
could choose \(p(N)+1\) orthonormal null vectors. Simultaneous
normal-slice extraction and the polarized norm identity in
Lemma~\ref{mass:slice-compactness} would give an orthonormal list
of that length in the direct sum of the \(p(N)\) one-dimensional
normal ground spaces, a contradiction.

Conversely, Lemma~\ref{mass:vacuum-recovery} supplies \(p(N)\)
asymptotically orthonormal even vectors with charge norm tending
to zero. Equation~\eqref{mass:kernel-gap} changes each by \(o(1)\)
when it is projected onto \(\Ker D_h\). Their projected Gram
matrix tends to the identity, so these projected vectors are
linearly independent for all sufficiently small \(h\). Thus the
massive kernel has dimension exactly \(p(N)\) there, is even,
and has the stated uniform gap. Its graded index is consequently
\(p(N)\). Index constancy in Lemma~\ref{mass:common-domain}
proves \eqref{eq:massive-index} for every \(h>0\).
\end{proof}

The massive signed count is now fixed for every positive coupling. To
compare it with the undeformed kernel, we next analyze large-coupling
states near separated block centers.
